\renewcommand{\thechapter}{1}
\chapter{Introduction}
\label{chapter:10-Introduction}
Graph data structures and the algorithms that operate on them are foundational to computer science.
As the social media platforms of the early $21^{st}$ century have proliferated into the global consciousness, graphs have increasingly been used as powerful abstractions to represent interconnected data collected from the real world. 
Recent years have seen a `Cambrian explosion' of artificial intelligence (AI) applications, and the hardware accelerators used to improve their efficency~\cite{Regli2022b}.
There is a significant overlap between the problems that have traditionally been in the domain of graph-processing and what modern, neural approaches to machine learning are advancing.
We term the overlap of graph processing and AI and \gai.
For \gai to be useful in the real world it must be capable of processing \textit{massive}, \textit{heterogeneous}, \textit{dynamic} graphs in \textit{resource-constrained} environments. 
As a result, the efficiency and effectiveness of \gai is critical.

\textbf{Massive.} 
The largest available `real world' graph datasets are over $3.5$ billion vertices and $128$ billion edges~\cite{Dhulipala2021}, and the `huge' synthetic benchmark from Graph500 has over $4.3$ trillion vertices and over $70$ trillion edges\footnote{\url{https://graph500.org/?page_id=12\#sec-3_4}}.
In uncompressed formats using inefficient algorithms, these sizes are well beyond the capability of any single commercial machine's memory, and sequential processing is extremely unlikely to be completed within any reasonable timeframe.
As available data continues to increase in volume, the problems caused by massive datasets will only get worse.

\textbf{Heterogeneous.}
The `huge' Graph500 benchmark graph is a simple, unattributed, undirected, unweighted graph with a (mostly) regular structure.
Real-world graphs are only analyzed because they contain interesting information. 
The topological structure of vertices and edges holds huge amounts of latent information.
It is the difference in the shapes of graphs that provide much of their utility, while also introducing complications in computation because the structure of each graph is so varied.
The attributes of the vertices and edges hold huge amounts of explicitly encoded information.
In particular, as \textit{Knowledge Graphs} (KG) become increasingly common they introduce complications of large amounts of data being stored in addition to the topology of the graph and often yield complex structures like multi-graphs and hyper-edges that dramatically increase the complexity of graph processing. 
Effective \gai solutions need to handle the heterogeneity of graph shape and graph content without unacceptable performance degradation. 

\textbf{Dynamic.}
The velocity of big data is well established. 
In \gai a high velocity of data, combined with the need to quickly generate solutions drive the adoption of dynamic (streaming) graph architectures. 
Dynamic graphs mean that the cost of constructing the graph cannot be divorced from its analysis, meaning that representations with slow, expensive update procedures are unviable. 
Further, live-data links to a demand for live-answers, and algorithms that cannot run in near-real time are unsuitable for these applications. 
Dynamic graph processing is under-explored in traditional graph processing, and almost completely untouched in the context of \gai.
Effective \gai solutions must support dynamic graphs to remain relevant into the future. 

\textbf{Resource Constrained.}
Processing massive heterogeneous graphs quickly enough to be useful is an expensive undertaking. 
Budgets are limited, so simply adding more computing resources that increase capital and operating costs does not work economically, even for the few \gai tasks where it makes sense computationally. 
The value of streaming data decays over time, and the longer it takes to generate insights, the less impactful they are. 
Further, because of storage and other limitations, data is not perpetually available, so processing within a given time budget becomes an important challenge. 
Commercial computers hit practical limits in memory and compute capacity that drives memory and compute efficiency. 
Distributed computing as a solution then introduces the need for communication efficiency.
Finally, as the push to `move processing to where the data is' motivates deploying \gai at-the-edge, Size, Weight and Power (SWaP) budgets will constrain the processing model and so incremental, approximate and other resource-preserving approaches must be developed to achieve tasks in constrained environments.   

The remainder of Chapter~\ref{chapter:10-Introduction} will define \gai, explain its importance in the real world and provide an overview of the challenges in understanding and improving \gai



%\section{Intelligence and Security Context}
\label{section:11-Intelligence-Security-Context}

% War is fundamentally a human activity, long understood to have an enduring \textit{nature} derived from its human-centricity and variable \textit{character} reflective of the context in which it is fought.  
% It has long been characterized by the \textit{danger} of armed conflict, the \textit{uncertainty} induced by the fog of war, the \textit{friction} generated by operating in a contested environment and pure \textit{chance}~\cite{Clausewitz2003}.


Many intelligence and security problems can be addressed computationally... 


\subsection{Use Cases for Intelligence and Security Problems}
\label{section:IS-Use-Cases}

%Overview of use-cases here. 

\subsubsection{Link Analysis}
\label{section:Link-Analysis}
% Reachability / Shortest Path
A link analysis problem requires an intelligence analyst to determine if two entities, events or other phenomena are related. 
Frequently a question of social contact, an analyst will build a list of known associates for each entity in question and expand those lists outwards until identifying a connection between the two. 
The problem can be multidimensional, friends list on a social media platform may give one subset of connections, the call data logs for the entities may give another and the locations that emerge in their pattern of life may form a third. 
Traditionally approached with pictures and red twine, more recently applications like the Analyst's Notebook\footnote{\url{https://i2group.com/i2-analysts-notebook}} are used to create digital equivalents.
Asking \textit{``Does this target associate with that target''} is then a problem of following the twine to see where (if at all) the two are connected. 

\subsubsection{Kill Chain Analysis}
\label{section:Kill-Chain-Analysis}
% Reachability / Shortest Path (Physical, Cyber)
A kill chain analysis is a popular methodology for post hoc incident analysis. 
The general idea is that for an incident to occur, an ordered set of events must precede it to the conditions for the incident. 
The canonical example draws from the counter-terrorism literature of the early 2000s, which asserts that for a terrorist attack to occur an attacker must select the target, then reconnoiter it, select specific targets, surveil the target, rehearse the attack, conduct the attack and then exploit their successes and escape~\cite{Tradoc2007}.
By understanding the chain of events, intelligence analysts seek to identify indicators and warnings from `left of bang' that may indicate an attack is imminent. 
To prevent an attacker from exploiting their success, you would block their actions on the objective. 
To prevent the attack, you catch them in rehearsals. 
To stop them from rehearsing, you should catch them surveilling the target, and so on.
Kill chains have also been popularised in the cyber domain by Lockheed Martin~\cite{Hutchins2011}.
It follows the same idea that the goal of a cyber threat intelligence analyst should be to detect adversary activity long before they can act on their objective. 
The challenge facing the analysts is how to identify when these events, or strings of events are occurring as part of that kill chain, rather than being indicative of random noise.

\subsubsection{Call Chaining Analysis}
\label{section:Call-Chaining}
% Graph Search

\subsubsection{Intelligence Preparation of the Battlespace (Population Segmentation)}
\label{section:IPB}
% Community Detection / Graph Segmentation

\subsubsection{Entity Resolution}
\label{section:Entity-Resolution}
% Entity Resolution (Time Constrained)

\subsubsection{Event and Entity Deduplication}
\label{section:Event-Entity-Deduplication}
% Entity Resolution

\subsubsection{Co-Traveler Analysis}
\label{section:Co-Traveler-Analysis}
% Contact Tracing, Drug Interdiction etc
% Algorithm??

\subsubsection{Financial Analysis}
\label{section:Financial-Analysis}
% Subgraph Matching

\subsubsection{Link Prediction}
\label{section:Link-Prediction}
% GNN Missing Link Prediction

\subsubsection{Intrusion Detection}
\label{section:Intrusion-Detection}

\subsubsection{Course of Action Prediction}
\label{section:Course-Of-Action-Prediction}

%
% Civy world - Recommend a friend on a platform (i.e. id who is most suited to be the bridge node -- MIT research?)

\subsubsection{Object Detection and Image Classification}
\label{section:Object-Detection-Image-Classification}
% Does Graph Accelerator improve classical NN speeds? 


\section{Graph-Based Artificial Intelligence}
\label{section:11-graph-based-ai}

\subsection{A Taxonomy of \gai}
Graph-Based Artificial Intelligence (\gai) is a sub-field of computer science dedicated to developing models and algorithms to support learning and reasoning over information represented in graph data structures.
Graphs are powerful abstractions of the real world that capture how entities relate to each other across a variety of dimensions, including social, spatial, semantic and logical. 
%
While graphs can express a wealth of information about the world, applying AI to that information effectively is challenging.
%
Most \gai approaches focus on analyzing the graph topology, suitable for exploring search spaces in planning problems, and basic pattern matching. 
For problems that require reasoning over the attributes of the graph, or the combination of its topology and attributes, the field still lacks appropriate approaches for dynamic graphs, heterogeneous graphs and graphs at the exa-scale. 
Most recently, the focus has shifted to using graphs to improve other AI models, such as with the Retrieval Augmented Generation (RAG) paradigm, where symbolic AI reasoning over graphs grounds generated content in curated, traceable facts.

\gai problems can be broken down into several overarching "AI Problem" categories (\textit{Searching, Planning \& Scheduling; Structural Analysis; Reasoning; Graph Abstraction; Graph Prediction} and \textit{Trustworthy AI}) which are outlined in Table~\ref{table:taxonomy}. 
Each \gai problem lists several typical tasks within the \gai problem space. 
The \gai taxonomy highlights the breadth of problems that a framework for benchmarking \gai must cover and emphasizes how difficult it is to fairly gauge `good performance' across \gai tasks.

\begin{table*}[t]
    \begin{center}
    \small
        \begin{tabular}{lll} 
            Relation to AI     &  AI Problem    &       Example Workloads      \\%&       \textbf{Citation(s)}    \\
            \hline
            & & \\
            Formulated as graph problem & Search                   & BFS, DFS, UC, A*, MinMax              \\%& \cite{Russell2021}            \\
                                        & Constraint Satisfaction  & Forward Checking, Backtracking        \\%& \cite{Russell2021}            \\
                                        & Planning                 & GraphPlan                             \\%& \cite{Russell2021}            \\
                                        & Scheduling               & CriticalPath                          \\%& \cite{Russell2021}            \\
                                        & Inference                & Variable Elimination (Bayes Net)     \\% & \cite{Russell2021}            \\
                                        & Reinforcement Learning   & Q-Learning                           \\% & \cite{Russell2021}            \\
            & & \\
            Operates on graph           & Pattern Matching         & Spatial Subgraph Matching             \\%& \cite{Osullivan2023}          \\
                                        & Graph Clustering         & Community Detection                  \\% & \cite{Girvan2002}             \\
                                        & Graph Partitioning       & Modularity Optimization               \\%& \cite{Brandes2006}            \\
                                        & Influence Mapping        & PageRank, BetweennessCentrality       \\%& \cite{Page1998, Brandes2001}  \\
                                        & Representation Learning  & Node2Vec, DeepWalk                    \\%& \cite{Grover2016, Perozzi2014} \\
            & & \\
            Operates on knowledge graph & Knowledge Acquisition & KG Completion, Ontology Learning   \\%& \cite{Ji2021, Maedche2004}          \\
                                        & Knowledge Representation Learning & Complex Vector Embeddings \\%& \cite{Ji2021} \\ 
                                        & Logic-Based Reasoning             & Logical Inference, Ontological Inference \\%& \cite{Chen2020} \\
                                        & Representation Based Reasoning    & Tensor Factorization, Semantic Matching  \\%& \cite{Chen2020} \\
                                        & Neural Reasoning                  & Relation Inference, Entity Resolution \\%& \cite{Chen2020} \\
            & & \\
            Graph Neural Networks       & Recurrent GNN             & Time Series Prediction               \\% & \cite{Wu2020} \\
                                        & Convolutional GNN         & Syntactic GCN, Scene-Graph Generation \\%& \cite{Zhang2019, Chang2021} \\
                                        & Graph AutoEncoders        & Link Prediction                       \\%& \cite{Wu2020} \\
                                        & Spatio-Temporal GNN       & Traffic Flow Prediction               \\%& \cite{Wu2020} \\
            & & \\
            Explainability \&  Interpretability              & Post-Hoc Explanations & XGNN                         \\% & \cite{Yuan2020} \\
                                        & Interpretable-By-Design  & GNN-Subset                           \\% & \cite{Pfiffer2022}\\

                                            
        \end{tabular}
    \normalsize
    \caption{A brief taxonomy of Graph-Based AI Problems} 
    \label{table:taxonomy}
    \end{center}
\end{table*}

\subsection{Real-world uses of \gai}
\gai is used heavily in real-world problem-solving. 
The breadth of tasks means that the `real-world' use varies from routine use to preliminary experimental work. 
A recent survey of \textit{Intelligence and Security} (I\&S) professionals by the \textit{Applied Research Laboratory for Intelligence and Security} (ARLIS)~\cite{Regli2022a} shows just how deeply \gai permeates. 

We extend that survey with a literature review (Table~\ref{table:problem-survey}), finding that many interesting problems need to examine the relationships between data, as well as the data itself. 
Graphs are highly suited to representing the relationships between data.
The problems span geospatial, social, financial and multi-sensor fusion domains, with many of the specific problems easily generalizable beyond an I\&S context.
Most input data are \textit{graph-native}, with the remainder easily represented as graph problems.
Of note, almost all the input data are heterogeneous attributed multigraphs, and many of the problems are expected to receive continuous update streams. 
Graph heterogeneity and streaming are two known sources of complexity in \gai.

The problems are spread across four of the six domains identified in the \gai Taxonomy (Table~\ref{table:taxonomy}).
\textit{Graph abstraction} is implicitly present in the \textit{graph prediction} (a downstream ML model can be used in place of a GNN), and \textit{Trustworthy AI} is an ongoing goal for all organizations, and will support human interfaces to these \gai tasks in the back-end. 
All of the tasks are constrained by the size of the input data and the complexity of representation. 
Of particular note, temporal limitations reflect either the requirement for near-real-time monitoring, or limited time to access data.

Overall, the problem survey shows that \gai problems deeply proliferate in the real world. 
The complexities introduced by heterogeneous representations, temporal constraints, streaming data and the breadth of tasks further emphasize the importance of characterizing \gai performance and understanding how interactions between the representations, programs, and systems manifest and the tradeoffs that need to be made when selecting the optimal approach to solve the problem(s) facing an organization.
 


\begin{landscape}
    
    \begin{table*}[t]
        \begin{center}
        \small
            \begin{tabular}{lllll} 
                AI Problem                                  & I\&S Problem                        & Input Data                       &  \gai Approach                 &   Constraints         \\
                \hline
                \textbf{\ac{GPS}}                                & Link Analysis~\cite{Xu2004}         & Knowledge Graph                   & Reachability             & Heterogeneity   \\
                %
                                                            & Navigation                          & Spatial graph                     & Path Planning            & SWaP; Complexity      \\
                %
                        ~                                   &                    ~                &                 ~                 &             ~                 &          ~            \\
                %
                \textbf{\ac{GM}}                                 & HVT Identification                  & Multi-Source Graphs               & Influence Mapping        & Heterogeneity         \\
                %
                                                            & Co-Traveller Detection~\cite{Sawas2019,Tritsarolis2021,Regli2022a,Tang2014,Wang2016a}              & Spatio-Temporal Data               & Clustering            & Temporal   \\
                %
                                                            & POI Identification~\cite{Osullivan2023,Osullivan2023a}                  & Location Data                     & Spatial Pattern Matching & Complexity \\
                
                                                            & Demographic Segmentation~\cite{Army2023,Zachary1977,Sobolevsky2014}            & Communication Logs, Social Media  & Community Detection      & Temporal     \\
                %
                                                            & Counter-Threat Finance~\cite{Keene2014,Balaskas2018,Goldsmith2020,Michalak2011,Soltani2016}              & Transaction Records, Blockchain   & Subgraph Matching        & Temporal     \\
                %
                                ~                           &                   ~                 &                 ~                 &                 ~             &          ~            \\
                %
                \textbf{\ac{GR}}                          & Kill-Chain Analysis~\cite{Tradoc2007,Hutchins2011,Sadlek2022,Osullivan2015}                & Event Logs                       & Logical Inference  & Completeness    \\
                %
                                                            & All-Source Fusion~\cite{Army2019,Regli2022a,Szekely2015}                   & Streaming Knowledge Graph         & Entity Resoultion  & Temporal   \\
                %
                                ~                           &           ~                         &                 ~                 &                 ~             &             ~         \\
                %
                \textbf{\ac{GP}}                         & Anomaly Detection~\cite{Wu2021,Ren2023}                   & Streaming Knowledge Graph         & Classification    & Temporal   \\
                %
                                                            & Link Prediction~\cite{Granovetter1973,Carmody2022,Bangerter2022, Ahmadi2023}                     & Streaming Knowledge Graph         & Completion        & Temporal   \\
                %
    

            \end{tabular}
        \normalsize
        \caption{A Survey of \gai Problems in the Intelligence \& Security Domain} 
        \label{table:problem-survey}
        \end{center}
    \end{table*}
\end{landscape}

\osullikomment{Maturity Model}


\subsection{What is the challenge with graph AI?}
\subsection{What do we need to fix the problem?}
\subsection{Why don't current benchmarks suffice?}
\subsection{How will you fix the problem?}


% Many of the `classical' AI problems in \textit{searching, planning and scheduling} start by abstracting the complexity of the real world to a graph-based formulation before searching for solutions, planning actions, making inferences or training agents~\cite{Russell2021}.
% When the input is a graph, AI approaches focus on the graph topology or the information in the attributes of the graph. 
% Topological problems necessitate \textit{structural analysis} to generate some useful knowledge or transformation, like matching locations based on spatial similarity~\cite{Osullivan2023} or detecting latent structures in the graph~\cite{Girvan2002, Brandes2006}.
% Information problems involve \textit{reasoning} over the graph abstraction of some knowledge with programs like Cyc~\cite{Lenat1990}, FrameNet~\cite{Fillmore2003}, ConceptNet~\cite{Speer2017}, ATOMIC~\cite{Sap2019} and COMET~\cite{Bosselut2019} garnering significant attention, inching towards the idea of a general-purpose common-sense AI that has driven research for the last sixty-odd years~\cite{Mccarthy1959}.
% As machine learning grew in popularity, the desire to apply machine learning to graph data structures followed.
% The initial push developed methods for \textit{graph abstraction} to embed graphs in vector space for classical machine learning applications.
% Neural approaches have yielded techniques for generating \textit{predictions} on graphs directly, with Graph Neural Networks (GNN) for classification, link prediction and traffic analysis providing even more powerful uses of graphs. 
% Finally, the structure of graphs and the long history of their use in reasoning are gaining attention as a mechanism to improve \textit{AI Trustworthiness}.
% For example, Retrieval Augmented Generation (RAG) research is exploring how to use knowledge graphs to improve Large Language Model (LLM) output and the use of symbolic AI to ground the outputs of generative models in verified, provable facts is an emerging area of interest. 

%
\section{Graph Data Structures}
\label{section:12_Graph_Data_Structures}
\section{Graph Hardware Acceleration}
\label{section:13-Graph-Hardware-Acceleration}
\section{Benchmarking Graphs}
\label{section:14-Benchmarking-Graphs}
\section{Thesis Outline}
\label{section:15-Thesis-Outline}
The thesis continues with a review of the \gai taxonomy and key graph and computer architecture concepts in Chapter~\ref{chapter:20-Preliminaries} to identify what makes optimizing \gai approaches hard and why measuring relative performance is difficult. 
Then, Chapter~\ref{chapter:30-Literature-Review} is a detailed review of current \gai benchmarking and workload analysis efforts to characterize the current gaps in benchmark coverage of \gai tasks, datasets, and metrics and identify what causes performance to degrade in \gai systems. 
Chapter~\ref{chapter:40-Graph-AI-Benchmarking} introduces the \gai benchmark frameworks and outlines our methodology for developing \gai problems to constitute the benchmark framework. 
We use the population segmentation problem from the Intelligence and Security domain to demonstrate how our development methodology is applied, describe the implementation of the \gai benchmark prototype, and outline the findings of our initial experiments in chapter~\ref{chapter:50-Community-Detection-Case-Study}.
Finally, chapter~\ref{chapter:60-Conclusion} concludes after identifying specific directions for future work. 